% !TEX program = xelatex
% xelatex -> bibtex -> xelatex -> xelatex
\documentclass{base}
\usepackage[utf8]{inputenc}

\title{blcu-thesis(no after-sale service)}
\author{}
\date{February 2022}

% 进行个人信息设置
\title{这是一个很长的这是一个很长的这是一个很长的毕业论文题目}
\author{某~某~某}    % 作者姓名
\date{2\,0\,2\,1~年~3~月}  % 完成日期
\college{学~院~名~称}
\major{专~业~名~称}  % 专业名称
\study{专~业~方~向~名~称}
\stunum{123000678}  % 学号
\instructor{某~某~某~~~教~授}  % 导师姓名

% 添加自己要用的其他宏包
\usepackage[numbers]{natbib}
\usepackage{subfig}
\usepackage{ctex}

% --- 证明结束黑框 ----
% \renewcommand{\qedsymbol}{$\blacksquare$}
% --- 设置字体 -----
\usepackage{newtxtext}  % for text fonts
\usepackage{xeCJK}
\setCJKmainfont{SimSun.ttf}
\setmainfont{Times New Roman}
% --- 设置数学字体 -----
% \usepackage{newtxmath}
% \usepackage{mathptmx}

% --- 直接插入 pdf 文件 ----
% \usepackage{pdfpages}

% --- 自定义命令 -----
\newcommand{\CC}{\ensuremath{\mathbb{C}}}
\newcommand{\RR}{\ensuremath{\mathbb{R}}}
\newcommand{\A}{\mathcal{A}}
\newcommand{\ii}{\bm{\mathrm{i}}\,}  % 虚部
\newcommand{\md}{\mathrm{d}\,}
\newcommand{\bA}{\boldsymbol{A}}
\newcommand{\red}[1]{\textcolor{red}{#1}}

\begin{document}

%%%%%%%%%%%%%%%%%%%%%%%%%%%%%%%%%%%%%%%%%%%%%%%%%%%
%\maketitle
%\makestatement

% \clearpage   % 结束上一页

% \begin{cnabstract}{关键词 1；关键词 2；关键词 3}

% 摘要内容摘要内容摘要内容摘要内容摘要内容摘要内容摘要内容摘要内容摘要内容摘要内容摘要内容摘要内容摘要内容摘要内容摘要内容摘要内容摘要内容摘要内容摘要内容摘要内容摘要内容摘要内容摘要内容摘要内容摘要内容.

% \end{cnabstract}

% \begin{enabstract}{Keyword 1;~ Keyword 2;~ Keyword 3}

% This is abstract. 

% \end{enabstract}
%%%%%%%%%%%%%%%%%%%%%%%%%%%%%%%%%%%%%%%%%%%%%%%%%%%	

\maketoc[pagenumtoc]
\clearpage 
\pagenumbering{arabic} 

\mainmatter
\chapter{绪论}
\section{选题背景}
\section{选题背景}
\section{选题背景}
\section{选题背景}
\chapter{绪论}
\section{选题背景}
\section{选题背景}
\section{选题背景}
\section{选题背景}
\chapter{绪论}
\section{选题背景}
\section{选题背景}
\section{选题背景}
\section{选题背景}
\chapter{绪论}
\section{选题背景}
\section{选题背景}
\section{选题背景}
\section{选题背景}

\begin{figure}[h]
\centering
  \includegraphics[width=0.45\linewidth]{paimon.png} 
  \caption{图片测试。}
  \label{fig:image}
  \vspace{-0.1cm}
\end{figure}

\begin{equation}\label{eqn:asr01}
    \begin{aligned}
    W^*&=\mathrm{argmax}_WP(W|O) \\
    &=\mathrm{argmax}_W\frac{P(O|W)P(W)}{P(O)} \\
    &\propto\mathrm{argmax}P(O|W)P(W) 
    \end{aligned}
\end{equation}

\section{参考文献}

参考文献

\backmatter

\bibliographystyle{shnuthesis-numeric}
\bibliography{mybib}
\begin{researchpage}
\hangindent 1.4em
\noindent
[1] {{\bf Author 1} and Author 2},  The name of the published article 1, {\bf  Name of Journal}, 2020, 12(34):1001-1020.

\hangindent 1.4em
\noindent
[2] {{\bf Author 1},  Author 2 and Author 3}, The name of the published article 2, submitted to
Journal of XXX.


\end{researchpage}

\begin{thankpage}
\chaptermark{致谢}
\setlength{\baselineskip}{24pt}

感谢老师

\end{thankpage}


%%%%%%%%%%%%%%%%%%%%%%%% 附录  %%%%%%%%%%%%%%%%%%%%%%%%%%%%%%%%%%%%%%%

% 添加附录, 如不需要可以注释掉
\appendix

%\setcounter{page}{1} % 如果需要可以自行重置页码
\renewcommand{\chaptermark}[1]{\markboth{#1}{}}

\chapter{附录 A ~ 这是第一个附录}
\renewcommand{\thesection}{A.\arabic{section}}
\section{附录A的小节}

这里是附录环境, 手动设置了 chapter 和 section 的样式, 并且加入到了目录.

\chapter{附录 B ~ 这是第二个附录}
\renewcommand{\thesection}{B.\arabic{section}}
\section{附录B的小节}

这里是附录环境, 手动设置了 chapter 和 section 的样式, 并且加入到了目录.


\end{document}